\chapter{Probabilidad sobre espacios numerables}\label{ch:b2:proba}

Los tres capítulos finales desarrollan la teoría moderna de la
probabilidad: \omterm{def:b2:proba:space}{medidas de probabilidad} sobre \omterm{def:b2:proba:space}{espacios muestrales}
\omterm{def:b2:structures:countable}{numerables}, variables aleatorias discretas y
\omterm{ex:b2:powerseries:fibonacci}{funciones generatrices}. La teoría finita del volumen de secundaria
adquiere toda su infraestructura: la $\sigma$-aditividad sustituye a
la aditividad finita, y la maquinaria de las \omterm{def:b2:series:summable}{familias sumables} del~\cref{ch:b2:series} es exactamente lo que hace manejables los
\omterm{def:b2:proba:space}{espacios muestrales} infinitos. Los resultados centrales de aquí son la
\omterm{def:b2:metric:continuity}{continuidad} de la probabilidad a lo largo de sucesiones monótonas de
\omterm{def:b2:proba:space}{sucesos} y el lema de Borel--Cantelli.

\section{Espacios de probabilidad}

\begin{definition}[Espacio de probabilidad numerable]\label{def:b2:proba:space}
Sea $\Omega$ un conjunto finito o \omterm{def:b2:structures:countable}{numerable} no vacío (el
\emph{espacio muestral}\index{espacio muestral}). Una \emph{medida de
probabilidad}\index{medida de probabilidad} sobre $\Omega$ es una
aplicación $\P$ del conjunto $\mathcal{P}(\Omega)$ de todos los
subconjuntos de $\Omega$ (los \emph{sucesos}\index{suceso}) en $[0,
1]$ tal que:
\begin{enumerate}
\item $\P(\Omega) = 1$;
\item ($\sigma$-aditividad) para toda sucesión $(A_n)_{n\in\N}$ de
sucesos disjuntos dos a dos,
\[
\P\Bigl(\,\bigcup_{n \in \N} A_n\Bigr)
= \sum_{n=0}^{\infty} \P(A_n) .
\]
\end{enumerate}
El par $(\Omega, \P)$ es un \emph{espacio de
probabilidad}\index{espacio de probabilidad} (\omterm{def:b2:structures:countable}{numerable}).
\end{definition}

\begin{remark}
Sobre un $\Omega$ \omterm{def:b2:structures:countable}{numerable} podemos tomar como \omterm{def:b2:proba:space}{sucesos} todos los
subconjuntos; sobre espacios no \omterm{def:b2:structures:countable}{numerables} (como los que hacen falta
para los modelos \omterm{def:b2:metric:continuity}{continuos} del tercer año) esto ya no es posible, y se
restringe $\P$ a una colección adecuada de \omterm{def:b2:proba:space}{sucesos}, una
\emph{$\sigma$-álgebra}. Todas las fórmulas de este capítulo
sobreviven a esa generalización palabra por palabra.
\end{remark}

\begin{proposition}[Reglas elementales]\label{prop:b2:proba:rules}
Para \omterm{def:b2:proba:space}{sucesos} $A, B$ y una \omterm{def:b2:proba:space}{medida de probabilidad} $\P$:
$\P(\emptyset) = 0$; $\P$ es finitamente aditiva; $\P(A^c) = 1 -
\P(A)$; si $A \subseteq B$, entonces $\P(A) \leq \P(B)$; y
\[
\P(A \cup B) = \P(A) + \P(B) - \P(A \cap B) .
\]
\end{proposition}

\begin{proof}
Aplicar la $\sigma$-aditividad a $A_0 = \Omega$, $A_n = \emptyset$
($n \geq 1$) da $1 = 1 + \sum_{n\geq1}\P(\emptyset)$, luego
$\P(\emptyset) = 0$; y rellenar una unión disjunta finita con
conjuntos vacíos da entonces la aditividad finita. Lo demás se sigue
como en el caso finito (volumen de secundaria): $1 = \P(A) + \P(A^c)$
a partir de $\Omega = A \sqcup A^c$; $\P(B) = \P(A) + \P(B \setminus
A) \geq \P(A)$ cuando $A \subseteq B$; y descomponiendo en tres
piezas disjuntas,
\begin{align*}
\P(A \cup B) &= \P(A \setminus B) + \P(B \setminus A) +
\P(A \cap B)\\
&= \bigl(\P(A) - \P(A\cap B)\bigr) + \bigl(\P(B) - \P(A\cap
B)\bigr) + \P(A \cap B),
\end{align*}
que es la inclusión-exclusión; la versión general con $n$ conjuntos
es el~\cref{exo:b2:proba:4}.
\end{proof}

\begin{proposition}[Distribuciones sobre un espacio
numerable]\label{prop:b2:proba:distribution}
Dar una \omterm{def:b2:proba:space}{medida de probabilidad} sobre un $\Omega = \{\omega_0,
\omega_1, \dots\}$ \omterm{def:b2:structures:countable}{numerable} equivale exactamente a dar pesos $p_i
= \P(\{\omega_i\}) \geq 0$ con $\sum_i p_i = 1$; y entonces, para
todo $A \subseteq \Omega$,
\[
\P(A) = \sum_{\omega \in A} \P(\{\omega\}) ,
\]
una subsuma (\omterm{def:b2:series:def}{absolutamente} convergente) de la familia $(p_i)$.
\end{proposition}

\begin{proof}
Dada $\P$, los conjuntos unitarios $\{\omega\}$, $\omega \in A$,
forman un recubrimiento disjunto \omterm{def:b2:structures:countable}{numerable} de $A$, de modo que la
$\sigma$-aditividad fuerza
\[
\P(A) = \sum_{\omega\in A}\P(\{\omega\}),
\]
una subsuma incondicional de la \omterm{def:b2:series:summable}{familia sumable} no negativa $(p_i)$;
reordenar es inofensivo precisamente porque los términos son no
negativos (\cref{ch:b2:series}); en particular, $\sum_ip_i =
\P(\Omega) = 1$. Recíprocamente, dados pesos no negativos de suma
total $1$, defínase $\P(A) = \sum_{\omega \in A}p_\omega$: la familia
es \omterm{def:b2:series:summable}{sumable}, y la $\sigma$-aditividad es exactamente el teorema de
sumación por paquetes del~\cref{ch:b2:series} aplicado a la partición
de $\bigcup A_n$ en los $A_n$.
\end{proof}

\begin{example}[Modelo geométrico: esperando la primera
cara]\label{ex:b2:proba:geometric}
Lánzese repetidamente una moneda con probabilidad de cara $p \in
\intoo{0}{1}$, y sea $\Omega = \N^* \cup \{\infty\}$ el registro del
rango de la primera cara. Los pesos naturales son
\[
\P(\{k\}) = (1 - p)^{k-1}p
\quad (k \in \N^*),
\qquad
\P(\{\infty\}) = 0 ,
\]
una \omterm{def:b2:proba:space}{medida de probabilidad}, pues $\sum_{k\geq1}(1-p)^{k-1}p =
\frac{p}{1 - (1-p)} = 1$: con probabilidad $1$ el juego termina;
pero el \omterm{def:b2:proba:space}{espacio muestral} debe seguir conteniendo la posibilidad de que
no lo haga. La aditividad \omterm{def:b2:structures:countable}{numerable} es lo que permite afirmar
$\P(\text{el juego termina}) = \sum_k \P(\{k\})$.
\end{example}

\begin{theorem}[Continuidad monótona]\label{thm:b2:proba:continuity}
Sea $(A_n)$ una sucesión de \omterm{def:b2:proba:space}{sucesos}.
\begin{enumerate}
\item Si $A_n \subseteq A_{n+1}$ para todo $n$ (sucesión
\emph{creciente}), entonces $\P\bigl(\bigcup_n A_n\bigr) =
\lim_{n\to\infty} \P(A_n)$.
\item Si $A_n \supseteq A_{n+1}$ para todo $n$ (sucesión
\emph{decreciente}), entonces $\P\bigl(\bigcap_n A_n\bigr) =
\lim_{n\to\infty} \P(A_n)$.
\end{enumerate}
\end{theorem}

\begin{proof}
\emph{1.} Disjuntícese: sean $B_0 = A_0$ y $B_n = A_n \setminus
A_{n-1}$. Los $B_n$ son disjuntos dos a dos con $\bigcup_{k \leq n}
B_k = A_n$ y $\bigcup_n B_n = \bigcup_n A_n$. Por la
$\sigma$-aditividad y la aditividad finita,
\[
\P\Bigl(\bigcup_n A_n\Bigr)
= \sum_{n=0}^\infty \P(B_n)
= \lim_{N\to\infty}\sum_{n=0}^N \P(B_n)
= \lim_{N\to\infty}\P(A_N) .
\]
\emph{2.} Pásese a los complementarios: $(A_n^c)$ es creciente con
unión $\bigl(\bigcap A_n\bigr)^c$; aplíquese la parte 1: $1 -
\P(\bigcap A_n) = \lim (1 - \P(A_n))$.
\end{proof}

\begin{corollary}[Subaditividad numerable]\label{cor:b2:proba:subadd}
Para toda sucesión de \omterm{def:b2:proba:space}{sucesos}, $\P\bigl(\bigcup_n A_n\bigr) \leq
\sum_{n=0}^\infty \P(A_n)$.
\end{corollary}

\begin{proof}
La subaditividad finita $\P(A_0 \cup \dots \cup A_N) \leq \sum_0^N
\P(A_n)$ se sigue de la inclusión-exclusión por inducción (o de la
aditividad sobre los $B_n \subseteq A_n$ disjuntizados). Hágase $N
\to \infty$: el miembro izquierdo converge a $\P(\bigcup_n A_n)$ por
la \omterm{def:b2:metric:continuity}{continuidad} monótona aplicada a la sucesión creciente $C_N = A_0
\cup \dots \cup A_N$.
\end{proof}

\begin{example}[La cota de la unión: burda pero
indestructible]\label{ex:b2:proba:unionbound}
La subaditividad con un número finito de \omterm{def:b2:proba:space}{sucesos} ---la \emph{cota de
la unión}--- cambia precisión por universalidad. Para el problema del
cumpleaños con $23$ personas, acotar la probabilidad de colisión por
la suma sobre parejas da
\[
\P(\text{colisión}) \leq \binom{23}2\cdot\frac1{365}
= \frac{253}{365} \approx 0.693 ,
\]
frente al valor verdadero $0.507$: se pasa por mucho, porque las
colisiones se solapan. Y sin embargo, la cota \emph{no} necesita
\omterm{def:b2:proba:independence}{independencia}, ni ley conjunta, ni nada más que las probabilidades
de las parejas; por eso, en el problema de fin de semana y a lo largo
del~\cref{ch:b2:randomvar}, la cota de la unión es la primera
herramienta que se saca: cuando resulta ser pequeña, el asunto queda
zanjado sin más modelización.
\end{example}

\begin{example}[Un seis llega, tarde o temprano]\label{ex:b2:proba:sixeventually}
Lánzese un dado equilibrado para siempre y sea $B_n = {}$``al menos
un seis entre las $n$ primeras tiradas'', una sucesión creciente de
\omterm{def:b2:proba:space}{sucesos} con $\P(B_n) = 1 - (5/6)^n$. La \omterm{def:b2:metric:continuity}{continuidad} monótona da
\[
\P(\text{aparece un seis tarde o temprano})
= \P\Bigl(\bigcup_nB_n\Bigr)
= \lim_n\bigl(1 - (5/6)^n\bigr) = 1 .
\]
Lo importante no es el límite (obvio), sino el paso lógico: ``tarde o
temprano'' es un \omterm{def:b2:proba:space}{suceso} acerca de \emph{infinitas} tiradas, fuera del
alcance de la aditividad finita, y la \omterm{def:b2:metric:continuity}{continuidad} monótona ---es
decir, la $\sigma$-aditividad--- es precisamente el axioma que le
asigna una probabilidad. Todo enunciado casi seguro del resto de este
libro pasa por esta misma puerta estrecha.
\end{example}

\section{Condicionamiento e independencia}

\begin{definition}[Probabilidad condicionada]\label{def:b2:proba:conditional}
Para \omterm{def:b2:proba:space}{sucesos} $A, B$ con $\P(B) > 0$, la \emph{probabilidad
condicionada}\index{probabilidad condicionada} de $A$ dado $B$ es
\[
\P(A \mid B) = \frac{\P(A \cap B)}{\P(B)} .
\]
La aplicación $A \mapsto \P(A \mid B)$ es a su vez una
\omterm{def:b2:proba:space}{medida de probabilidad} sobre $\Omega$.
\end{definition}

\begin{remark}
Que $A \mapsto \pcond BA$ sea de nuevo una \omterm{def:b2:proba:space}{medida de probabilidad}
merece un momento: $\pcond B\Omega = 1$ y la $\sigma$-aditividad
pasan por el cociente porque la intersección con $B$ respeta las
uniones disjuntas. La consecuencia práctica: toda identidad de este
capítulo ---inclusión-exclusión, \omterm{def:b2:metric:continuity}{continuidad} monótona,
Borel--Cantelli--- puede aplicarse \emph{después} de condicionar, sin
demostraciones nuevas. Los probabilistas ``trabajan bajo $\pcond
B{\cdot}$'' constantemente por esta razón exacta.
\end{remark}

\begin{example}[Condicionar puede crear
uniformidad]\label{ex:b2:proba:sumseven}
Lánzense dos dados equilibrados y condiciónese a que la suma sea $7$:
para cada $k \in \intint16$,
\[
\pcond{\{S = 7\}}{X = k}
= \frac{\P(X = k,\ Y = 7 - k)}{\P(S = 7)}
= \frac{1/36}{6/36} = \frac16 :
\]
dada una suma de $7$, el primer dado es exactamente uniforme; $7$ es
el único total compatible con todas las caras, así que el
condicionamiento borra toda la información sobre $X$. Cualquier otro
total sesga la ley (dado $S = 4$, el primer dado es uniforme solo
sobre $\{1, 2, 3\}$). Calcular una ley condicionada consiste en
renormalizar los pesos conjuntos a lo largo del \omterm{def:b2:proba:space}{suceso}
condicionante, y en nada más.
\end{example}

\begin{example}[La segunda extracción es tan buena como la
primera]\label{ex:b2:proba:seconddraw}
Una urna contiene $3$ bolas blancas y $2$ negras; extráiganse dos sin
reemplazamiento. Todos coinciden en que $\P(W_1) = \frac35$; ¿cuánto
vale $\P(W_2)$? Probabilidad total a lo largo de la primera
extracción:
\[
\P(W_2) = \pcond{W_1}{W_2}\,\P(W_1) +
\pcond{B_1}{W_2}\,\P(B_1)
= \frac24\cdot\frac35 + \frac34\cdot\frac25
= \frac{12}{20} = \frac35 :
\]
exactamente $\P(W_1)$. No hacía falta ningún cálculo: por simetría,
toda bola tiene la misma probabilidad de ser la segunda extraída, así
que la segunda extracción ---\emph{incondicionalmente}--- tiene la
misma ley que la primera. Condicionar al primer resultado cambia las
probabilidades; no conocerlo, no. Este argumento de
intercambiabilidad regresa en el capítulo siguiente para el muestreo
sin reemplazamiento, donde da la media hipergeométrica $np$ sin
ninguna identidad binomial.
\end{example}

\begin{theorem}[Probabilidades compuestas, probabilidad total,
Bayes]\label{thm:b2:proba:bayes}
\leavevmode
\begin{enumerate}
\item (Regla de la cadena) Si $\P(A_1 \cap \dots \cap A_{n-1}) > 0$,
\[
\P(A_1 \cap \dots \cap A_n)
= \P(A_1)\,\P(A_2 \mid A_1)\cdots
\P(A_n \mid A_1 \cap \dots \cap A_{n-1}) .
\]
\item (Probabilidad total) Si $(B_i)_{i \in I}$ es una partición
finita o \omterm{def:b2:structures:countable}{numerable} de $\Omega$ con $\P(B_i) > 0$, entonces, para
todo \omterm{def:b2:proba:space}{suceso} $A$:
\[
\P(A) = \sum_{i \in I} \P(A \mid B_i)\,\P(B_i) .
\]
\item (Bayes) Bajo las mismas hipótesis, si además $\P(A) > 0$:
\[
\P(B_j \mid A)
= \frac{\P(A \mid B_j)\,\P(B_j)}
       {\sum_{i \in I} \P(A \mid B_i)\,\P(B_i)} .
\]
\end{enumerate}
\end{theorem}

\begin{proof}
\emph{1.} Escríbase cada \omterm{def:b2:proba:conditional}{probabilidad condicionada} como cociente: el
segundo miembro es
\[
\P(A_1)\cdot\frac{\P(A_1 \cap A_2)}{\P(A_1)}\cdot
\frac{\P(A_1 \cap A_2 \cap A_3)}{\P(A_1 \cap A_2)}\cdots
\frac{\P(A_1 \cap \dots \cap A_n)}{\P(A_1 \cap \dots \cap
A_{n-1})},
\]
un producto telescópico: cada denominador cancela el numerador
anterior y deja $\P(A_1 \cap \dots \cap A_n)$. Todos los
denominadores son $\geq \P(A_1 \cap \dots \cap A_{n-1}) > 0$ por
monotonía, así que nada se anula. (La hipótesis guarda exactamente
esto: condicionar a un \omterm{def:b2:proba:space}{suceso} de probabilidad cero no está
definido.) \emph{2.} Los conjuntos $A \cap B_i$ son disjuntos dos a
dos y de unión $A$; aplíquese la ($\sigma$-)aditividad y la
definición de condicionamiento. \emph{3.} Ambos miembros de $\P(B_j
\mid A)\P(A) = \P(A \mid B_j)\P(B_j)$ valen $\P(A \cap B_j)$;
divídase por $\P(A)$ y desarróllese $\P(A)$ por probabilidad total.
\end{proof}

\begin{example}[La colisión de cumpleaños, por la regla de la
cadena]\label{ex:b2:proba:birthday}
Con $n$ personas cuyos cumpleaños son \omterm{def:b2:proba:independence}{independientes} y uniformes sobre
$365$ días, sea $D_n = {}$``los $n$ cumpleaños son todos distintos''.
Condicionando persona a persona (regla de la cadena):
\[
\P(D_n) = \prod_{k=1}^{n-1}\Bigl(1 - \frac{k}{365}\Bigr),
\]
pues cada nueva persona ha de evitar los $k$ días ya ocupados. Para
$n = 23$: $\P(D_{23}) \approx 0.493$; un cumpleaños compartido es ya
más probable que improbable. La heurística que explica la pequeñez de
$23$: tomando logaritmos, $-\ln \P(D_n) \approx \sum_{k<n}\frac
k{365} = \frac{\binom n2}{365}$, y $\binom{23}2 = 253$ da $253/365
\approx 0.693 \approx \ln 2$. Lo que importa es el número de
\emph{parejas}, que crece cuadráticamente: los problemas de colisión
viven en la escala $n \sim \sqrt{365}$, no en $n \sim 365$; la
paradoja del cumpleaños es una raíz cuadrada disfrazada.
\end{example}

\begin{example}[Monty Hall, por Bayes]\label{ex:b2:proba:montyhall}
Un premio se esconde tras una de tres puertas, uniformemente. Eliges
la puerta $1$; el presentador, que sabe dónde está el premio, abre
una de las otras puertas, siempre vacía (eligiendo uniformemente
cuando tiene opción), digamos la puerta $3$. Sean $B_i = {}$``el
premio está tras la puerta $i$'' y $A = {}$``el presentador abre la
puerta $3$''. Entonces $\pcond{B_1}{A} = \frac12$, $\pcond{B_2}{A} =
1$ y $\pcond{B_3}{A} = 0$, luego, por Bayes
(\cref{thm:b2:proba:bayes}),
\[
\P(B_2 \mid A)
= \frac{1\cdot\frac13}
{\frac12\cdot\frac13 + 1\cdot\frac13 + 0\cdot\frac13}
= \frac23 :
\]
cambiar de puerta gana dos de cada tres veces. El cálculo localiza
exactamente la confusión popular: el movimiento del presentador es
\emph{informativo} (no podría abrir la puerta $2$ si el premio
estuviera allí), y la fórmula de Bayes es el instrumento contable que
convierte esa asimetría en el $\frac23$. Condicionar a ``lo que se ha
visto'' en vez de a ``lo que es cierto'' es todo el arte de la
fórmula.
\end{example}

\begin{example}[Las dos apuestas del caballero de
Méré]\label{ex:b2:proba:demere}
Dos apuestas del siglo XVII, zanjadas por la \omterm{def:b2:proba:independence}{independencia}. Apuesta
uno: al menos un seis en $4$ tiradas de un dado,
\[
\P = 1 - \Bigl(\frac56\Bigr)^{\!4} \approx 0.518 > \frac12 .
\]
Apuesta dos: al menos un doble seis en $24$ tiradas de dos dados,
\[
\P = 1 - \Bigl(\frac{35}{36}\Bigr)^{\!24} \approx 0.491 <
\frac12 .
\]
De Méré razonaba que $24$ tiradas con probabilidad $\frac1{36}$
deberían igualar a $4$ tiradas con probabilidad $\frac16$ (misma
razón $\frac{24}{36} = \frac46$); el fallo de esa proporcionalidad
---las probabilidades de uniones no escalan linealmente--- fue, se
dice, lo que motivó su carta a Pascal y, con ella, el nacimiento de
la teoría de la probabilidad. La comparación correcta pasa por los
logaritmos: $n$ ensayos con probabilidad $p$ tienen al menos un éxito
con probabilidad $1 - (1-p)^n \approx 1 - \eu^{-np}$, de modo que el
invariante honesto es $np$: aquí, $4\cdot\frac16 = \frac23$ frente a
$24\cdot\frac1{36} = \frac23$; ¡iguales! Las dos apuestas difieren
solo en el segundo orden en $p$, y justo lo bastante para mover una
al otro lado de la línea del cincuenta por ciento: las probabilidades
pequeñas son un terreno en el que la intuición necesita la
exponencial, no la regla.
\end{example}

\begin{remark}[Falacias frecuentes del condicionamiento]
Tres confusiones recurrentes, todas visibles en los ejemplos
anteriores. (i) \emph{Inversión}: $\pcond BA$ y $\pcond AB$ difieren
en el factor $\P(A)/\P(B)$; un test con un $99\,\%$ de acierto sobre
los enfermos puede dejar a un paciente positivo casi con certeza sano
cuando la enfermedad es rara (\cref{exo:b2:proba:3}); citar
$\pcond{\text{enfermo}}{\text{positivo}}$ donde se quiere decir
$\pcond{\text{positivo}}{\text{enfermo}}$ es la falacia de la tasa
base. (ii) \emph{Condicionar al \omterm{def:b2:proba:space}{suceso} equivocado}: en Monty Hall, el
\omterm{def:b2:proba:space}{suceso} condicionante correcto es ``el presentador abrió la puerta
$3$'', no ``el premio no está tras la puerta $3$''; los dos llevan
información distinta, y todo el $\frac23$ depende de esa diferencia.
(iii) \emph{Disjuntos frente a \omterm{def:b2:proba:independence}{independientes}}: unos \omterm{def:b2:proba:space}{sucesos} disjuntos
de probabilidad positiva nunca son \omterm{def:b2:proba:independence}{independientes} ($\P(A\cap B) = 0
\neq \P(A)\P(B)$); la \omterm{def:b2:proba:independence}{independencia} es compatibilidad de información,
no ausencia de solapamiento.
\end{remark}

\begin{definition}[Independencia]\label{def:b2:proba:independence}
Dos \omterm{def:b2:proba:space}{sucesos} $A$ y $B$ son \emph{independientes}\index{sucesos
independientes} si $\P(A \cap B) = \P(A)\P(B)$. Una familia $(A_i)_{i
\in I}$ de \omterm{def:b2:proba:space}{sucesos} es \emph{(mutuamente) independiente} si para todo
subconjunto finito $J \subseteq I$,
\[
\P\Bigl(\bigcap_{i \in J} A_i\Bigr)
= \prod_{i \in J} \P(A_i) .
\]
\end{definition}

\begin{remark}
La \omterm{def:b2:proba:independence}{independencia} mutua es estrictamente más fuerte que la
\omterm{def:b2:proba:independence}{independencia} dos a dos: con dos lanzamientos de moneda equilibrada,
los \omterm{def:b2:proba:space}{sucesos} ``la primera es cara'', ``la segunda es cara'' y ``ambas
coinciden'' son \omterm{def:b2:proba:independence}{independientes} dos a dos (cada pareja tiene
probabilidad de intersección $\frac14 = \frac12\cdot\frac12$) y, sin
embargo, la intersección triple tiene probabilidad $\frac14 \neq
\frac18$. Nótese también que si $A, B$ son \omterm{def:b2:proba:independence}{independientes}, también lo
son $A, B^c$ (calcúlese: $\P(A \cap B^c) = \P(A) - \P(A\cap B) =
\P(A)(1 - \P(B))$), y por tanto también $A^c, B^c$.
\end{remark}

\begin{example}[Independencia leída en una estructura de
producto]\label{ex:b2:proba:productindep}
Lánzense dos dados equilibrados: $\Omega = \intint16^2$ con pesos
uniformes. Sean $A = {}$``el primer dado es par'' y $B = {}$``el
segundo dado es al menos $5$''. Contando: $\abs A = 3\cdot6 = 18$,
$\abs B = 6\cdot2 = 12$, $\abs{A\cap B} = 3\cdot2 = 6$, luego
\[
\P(A\cap B) = \frac6{36} = \frac{18}{36}\cdot\frac{12}{36} =
\P(A)\,\P(B) :
\]
\omterm{def:b2:proba:independence}{independientes}, y el mecanismo está a la vista: $A$ restringe solo la
primera coordenada, $B$ solo la segunda, y la medida uniforme sobre
un conjunto producto hace que los recuentos por coordenadas se
multipliquen. Toda afirmación del tipo ``los \omterm{def:b2:proba:space}{sucesos} que dependen de
grupos disjuntos de lanzamientos son \omterm{def:b2:proba:independence}{independientes}'' (usada
masivamente en el problema de fin de semana) es este cálculo, con más
índices.
\end{example}

\begin{example}[Análisis del primer paso]\label{ex:b2:proba:firststep}
Para el modelo geométrico del~\cref{ex:b2:proba:geometric}, ¿cuál es
la probabilidad $u$ de que la primera cara caiga en un rango
\emph{par}? Condiciónese al primer lanzamiento: con probabilidad $p$
el rango es $1$ (impar); con probabilidad $q = 1 - p$ el juego
recomienza con todas las paridades invertidas, de modo que
\[
u = p\cdot0 + q\,(1 - u)
\qquad\Longrightarrow\qquad
u = \frac{q}{1 + q} .
\]
Una línea, sin series; y concuerda con la sumación directa del~\cref{exo:b2:proba:9}, que da $1 - u = \frac1{1+q}$. Esta técnica del
``primer paso'' (condicionar al primer experimento y reconocer una
copia desplazada del problema) es la forma probabilista de una
recursión, y es el motor que hay detrás de las ecuaciones de duración
del juego del~\cref{exo:b2:proba:6} y de los cálculos de primer paso
del problema de fin de semana.
\end{example}

\section{El lema de Borel--Cantelli}

\begin{definition}[Límite superior de sucesos]\label{def:b2:proba:limsup}
Para una sucesión $(A_n)$ de \omterm{def:b2:proba:space}{sucesos}, el \omterm{def:b2:proba:space}{suceso}
\[
\limsup_n A_n
= \bigcap_{N=0}^{\infty}\ \bigcup_{n \geq N} A_n
= \{\omega \in \Omega : \omega \in A_n
\text{ para infinitos } n\}
\]
es el \omterm{def:b2:proba:space}{suceso} ``$A_n$ ocurre infinitas veces''.
\end{definition}

\begin{example}[Traducir ``infinitas veces'' y ``a partir de cierto
punto'']\label{ex:b2:proba:limsuplang}
El complementario de $\limsup_nA_n$ es, por De Morgan,
\[
\Bigl(\bigcap_N\bigcup_{n\geq N}A_n\Bigr)^{\!c}
= \bigcup_N\bigcap_{n\geq N}A_n^c
= \{\omega : \omega \notin A_n \text{ para todo }n\text{
grande}\},
\]
el \omterm{def:b2:proba:space}{suceso} ``\emph{a partir de cierto punto}, $A_n$ falla'' (escrito
$\liminf_nA_n^c$). Así pues, ``$A_n$ infinitas veces'' y ``$A_n^c$ a
partir de cierto punto'' son complementarios; tener claro este
diccionario evita la mayoría de los accidentes con cuantificadores.
Traducciones de muestra para el lanzamiento de moneda: ``infinitas
caras'' es $\limsup\{X_n = H\}$; ``solo un número finito de rachas de
$100$ caras'' es el complementario de un límite superior; ``la
frecuencia acumulada converge a $\frac12$'' es
$\bigcap_j\bigcup_N\bigcap_{n\geq N}\{\abs{\widehat p_n -
\tfrac12} < \tfrac1j\}$; operaciones \omterm{def:b2:structures:countable}{numerables} en todos los casos,
así que todos ellos son \omterm{def:b2:proba:space}{sucesos} honestos.
\end{example}

\begin{theorem}[Borel--Cantelli]\label{thm:b2:proba:borelcantelli}
\leavevmode
\begin{enumerate}
\item Si $\sum_{n} \P(A_n) < \infty$, entonces
$\P\bigl(\limsup_n A_n\bigr) = 0$.
\item Si los \omterm{def:b2:proba:space}{sucesos} $A_n$ son \omterm{def:b2:proba:independence}{independientes} y $\sum_n \P(A_n) =
\infty$, entonces $\P\bigl(\limsup_n A_n\bigr) = 1$.
\end{enumerate}
\end{theorem}

\begin{proof}
\emph{1.} Sea $C_N = \bigcup_{n \geq N}A_n$; la sucesión $(C_N)$ es
decreciente con intersección $\limsup A_n$ y, por la subaditividad
\omterm{def:b2:structures:countable}{numerable} (\cref{cor:b2:proba:subadd}),
\[
\P(C_N) \leq \sum_{n \geq N}\P(A_n)
\xrightarrow[N\to\infty]{} 0
\]
(cola de una serie convergente). La \omterm{def:b2:metric:continuity}{continuidad} monótona
(\cref{thm:b2:proba:continuity}) concluye: $\P(\limsup A_n) = \lim_N
\P(C_N) = 0$.

\emph{2.} Basta probar que $\P\bigl(\bigcup_{n\geq N}A_n\bigr) = 1$
para todo $N$: en efecto, si unos \omterm{def:b2:proba:space}{sucesos} $B_N$ tienen todos
probabilidad $1$, entonces
\[
\P\Bigl(\Bigl(\bigcap_NB_N\Bigr)^{\!c}\Bigr)
= \P\Bigl(\bigcup_NB_N^c\Bigr)
\leq \sum_N\P(B_N^c) = 0
\]
por la subaditividad \omterm{def:b2:structures:countable}{numerable} (\cref{cor:b2:proba:subadd}), de modo
que la intersección \omterm{def:b2:structures:countable}{numerable} $\limsup A_n =
\bigcap_N\bigcup_{n\geq N}A_n$ sigue teniendo probabilidad $1$.
Fíjese $N$ y considérese, para $M > N$, el complementario:
\[
\P\Bigl(\bigcap_{n=N}^{M} A_n^c\Bigr)
= \prod_{n=N}^{M}\bigl(1 - \P(A_n)\bigr)
\leq \prod_{n=N}^{M} e^{-\P(A_n)}
= \exp\Bigl(-\sum_{n=N}^M \P(A_n)\Bigr) ,
\]
usando la \omterm{def:b2:proba:independence}{independencia} de los complementarios y la cota de convexidad
$1 - x \leq e^{-x}$. Cuando $M \to \infty$, el exponente tiende a
$-\infty$ por la divergencia de la serie, de modo que, por la
\omterm{def:b2:metric:continuity}{continuidad} monótona (sucesión decreciente), $\P\bigl(\bigcap_{n \geq
N}A_n^c\bigr) = 0$, es decir, $\P\bigl(\bigcup_{n \geq N}A_n\bigr) =
1$.
\end{proof}

\begin{example}[Rachas infinitas de caras]\label{ex:b2:proba:runs}
Lánzese una moneda equilibrada para siempre y sea $A_n$ el \omterm{def:b2:proba:space}{suceso}
``los lanzamientos $n, n+1, \dots, n + k - 1$ son todos caras'' (una
racha de $k$ caras que empieza en el instante $n$), con $k$ fijo. Los
\omterm{def:b2:proba:space}{sucesos} $A_{jk}$ ($j = 1, 2, \dots$), que dependen de bloques
disjuntos de lanzamientos, son \omterm{def:b2:proba:independence}{independientes}, cada uno de
probabilidad $2^{-k}$, y $\sum_j 2^{-k} = \infty$: por
Borel--Cantelli~2, con probabilidad $1$ hay infinitos bloques de solo
caras; \emph{todo} patrón fijo se repite infinitas veces, casi
seguramente. Recíprocamente, si dejamos crecer la \omterm{def:b2:curves:length}{longitud} de la
racha, $B_n = {}$``una racha de $2\log_2 n$ caras empieza en $n$''
tiene $\P(B_n) = n^{-2}$ \omterm{def:b2:series:summable}{sumable}, de modo que casi seguramente solo
empiezan un número finito de rachas tan largas: Borel--Cantelli
calibra con precisión \emph{cuán largas} son las rachas más largas.
\end{example}

\begin{example}[El mono infinito, cuantificado]
Un mono teclea letras \omterm{def:b2:proba:independence}{independientes} y uniformes de un alfabeto de
$26$ letras. Córtese el mecanoescrito en bloques disjuntos de cuatro
letras; los \omterm{def:b2:proba:space}{sucesos} $A_j = {}$``el bloque $j$ deletrea MATH'' son
\omterm{def:b2:proba:independence}{independientes} con $\P(A_j) = 26^{-4}$, y $\sum_j\P(A_j) = \infty$:
por Borel--Cantelli~2, el mono teclea MATH infinitas veces, casi
seguramente; y lo mismo vale para cualquier texto fijo de cualquier
\omterm{def:b2:curves:length}{longitud}, ajustando los bloques. La nota al pie cuantitativa desinfla
el milagro: $26^4 = 456\,976$, así que el primer MATH tarda unas
medio millón de pulsaciones de media, y una obra de Shakespeare de
$10^5$ caracteres espera del orden de $26^{10^5}$ bloques; ``casi
seguro'' es un enunciado sobre el horizonte $\infty$, no sobre ningún
horizonte que vaya a encontrarse un mono. Borel--Cantelli certifica
el límite; el tamaño de los sumandos cuenta la historia a escalas
humanas.
\end{example}

\begin{remark}
En el~\cref{ex:b2:proba:runs}, el \omterm{def:b2:proba:space}{espacio muestral} subyacente
(sucesiones infinitas de lanzamientos) es no \omterm{def:b2:structures:countable}{numerable}, de modo que,
en rigor, el ejemplo vive en el marco de la teoría de la medida del
tercer año; los \emph{cálculos}, sin embargo, solo usan las reglas
demostradas en este capítulo, aplicadas a \omterm{def:b2:proba:space}{sucesos} determinados por un
número finito de lanzamientos y a sus combinaciones \omterm{def:b2:structures:countable}{numerables}. Este
es el convenio estándar a este nivel: la teoría se enuncia sobre
espacios \omterm{def:b2:structures:countable}{numerables}, y los ejemplos de juegos infinitos se tratan con
las mismas herramientas.
\end{remark}

\begin{remark}[Perspectivas dentro de este volumen]
La maquinaria de este capítulo la consumen enteramente los dos
siguientes. Las indicatrices convierten los \omterm{def:b2:proba:space}{sucesos} en variables
aleatorias, y la $\sigma$-aditividad pasa a ser la \omterm{def:b2:series:summable}{sumabilidad} que
define la esperanza (\cref{ch:b2:randomvar}); Borel--Cantelli más una
cota \omterm{def:b2:series:summable}{sumable} de la cola es exactamente como se demuestra allí la ley
fuerte de los grandes números para monedas. En el~\cref{ch:b2:genfun}, la \omterm{def:b2:metric:continuity}{continuidad} monótona reaparece en el momento
decisivo: la probabilidad de extinción de un proceso de ramificación
se \emph{define} como el límite monótono $\lim\P(Z_n = 0)$, y la
ecuación de punto fijo que satisface se obtiene pasando al límite en
esa sucesión creciente; el teorema final del libro se sostiene sobre
el primer teorema de este capítulo.
\end{remark}

\begin{remark}[Método: tres vías hacia la probabilidad uno]
Los enunciados casi seguros se demuestran con tres palancas, en orden
creciente de potencia. \emph{\omterm{def:b2:metric:continuity}{Continuidad} monótona}: exhíbase el
\omterm{def:b2:proba:space}{suceso} como unión creciente (o intersección decreciente) de \omterm{def:b2:proba:space}{sucesos}
de horizonte finito con probabilidades calculables
(\cref{ex:b2:proba:sixeventually}). \emph{Uniones nulas}: una unión
\omterm{def:b2:structures:countable}{numerable} de \omterm{def:b2:proba:space}{sucesos} de probabilidad cero es nula (subaditividad
\omterm{def:b2:structures:countable}{numerable}), de modo que basta matar cada \omterm{def:b2:proba:space}{suceso} malo por separado;
así es como ``para todo $j$, a partir de cierto punto $\abs{\widehat
p_n - p} < 1/j$'' se ensambla en una convergencia.
\emph{Borel--Cantelli}: cuando el \omterm{def:b2:proba:space}{suceso} es un límite superior, súmense
las probabilidades; la convergencia lo mata (sin necesidad de
\omterm{def:b2:proba:independence}{independencia}) y la divergencia más la \omterm{def:b2:proba:independence}{independencia} lo certifican.
Elegir la palanca adecuada suele ser toda la demostración; el
problema de fin de semana usa las tres en un mismo argumento.
\end{remark}

\begin{remark}[Dónde se usa]
La \omterm{def:b2:metric:continuity}{continuidad} monótona y Borel--Cantelli son las dos palancas de
todo enunciado ``casi seguro'': impulsan la recurrencia del
\omterm{pb:b2:proba:1}{paseo aleatorio} del problema de fin de semana de este capítulo, el
lado casi seguro de la ley de los grandes números
(\cref{ch:b2:randomvar}) y el análisis de la extinción de los
procesos de ramificación (\cref{ch:b2:genfun}). El volumen del tercer
año reconstruye la teoría sobre $\sigma$-álgebras e integración de
Lebesgue, donde los \omterm{def:b2:proba:space}{espacios muestrales} no \omterm{def:b2:structures:countable}{numerables} usados aquí de
manera informal se vuelven plenamente rigurosos.
\end{remark}

\section{Ejercicios}

\begin{exercise}[$\star$]\label{exo:b2:proba:1}
Una urna contiene $n$ bolas numeradas. Las bolas se extraen una a una
sin reemplazamiento. Calcula la probabilidad de que la bola número
$1$ se extraiga antes que la bola número $2$. Generaliza: la
probabilidad de que la bola $1$ sea la primera extraída entre las
bolas $1, \dots, k$.
\end{exercise}

\begin{exercise}[$\star$]\label{exo:b2:proba:2}
Prueba que sobre $\Omega = \N^*$ los pesos $p_k = \frac{1}{k(k+1)}$
definen una \omterm{def:b2:proba:space}{medida de probabilidad}, y calcula $\P(2\N^*)$ (resultados
pares) como serie; prueba que vale $1 - \ln 2$. \emph{(Telescopa
$\frac{1}{2j(2j+1)} = \frac{1}{2j} - \frac{1}{2j+1}$ y usa la serie
armónica alternada, \cref{ch:b2:series}.)}
\end{exercise}

\begin{exercise}[$\star$]\label{exo:b2:proba:3}
(Falsos positivos) Una enfermedad afecta a una persona de cada
$10\,000$. Un test la detecta con probabilidad $0.99$ en los enfermos
y da un falso positivo con probabilidad $0.01$ en los sanos. Calcula
la probabilidad de estar enfermo dado un test positivo, y coméntala.
\end{exercise}

\begin{exercise}[$\star\star$]\label{exo:b2:proba:4}
Sean $A_1, \dots, A_n$ \omterm{def:b2:proba:space}{sucesos}. Demuestra la fórmula de
inclusión-exclusión
\[
\P\Bigl(\bigcup_{i=1}^n A_i\Bigr)
= \sum_{\emptyset \neq J \subseteq \{1,\dots,n\}}
(-1)^{\abs J + 1}\,\P\Bigl(\bigcap_{i \in J}A_i\Bigr)
\]
integrando la identidad $1 - \prod_{i=1}^n(1 - \mathbf{1}_{A_i}) =
\mathbf{1}_{\bigcup A_i}$ sobre $\Omega$ (es decir, sumando ponderado
por $\P(\{\omega\})$).
\end{exercise}

\begin{exercise}[$\star\star$]\label{exo:b2:proba:5}
(Problema de los emparejamientos, vía inclusión-exclusión) Se meten
$n$ cartas uniformemente al azar en $n$ sobres, una en cada uno.
Usando el~\cref{exo:b2:proba:4}, prueba que la probabilidad de que
\emph{ningún} emparejamiento sea correcto es $\sum_{k=0}^n
\frac{(-1)^k}{k!} \to e^{-1}$, y deduce la probabilidad de que haya
exactamente un acierto.
\end{exercise}

\begin{exercise}[$\star\star$]\label{exo:b2:proba:6}
Se lanza una moneda sesgada (probabilidad de cara $p \in
\intoo{0}{1}$) hasta que aparecen dos caras consecutivas. Sea $q_n$
la probabilidad de que el juego dure más de $n$ lanzamientos. Prueba,
condicionando al primer lanzamiento (o a los primeros), que $q_n =
(1-p)\,q_{n-1} + p(1-p)\,q_{n-2}$ para $n \geq 2$, y deduce que el
juego termina con probabilidad $1$. \emph{(Prueba que $q_n \to 0$
comparando con una sucesión geométrica: ambas raíces de la ecuación
característica tienen valor absoluto en $\intoo{0}{1}$.)}
\end{exercise}

\begin{exercise}[$\star\star\star$]\label{exo:b2:proba:7}
(Récords) Sórtese una sucesión infinita de ordenaciones uniformes
\omterm{def:b2:proba:independence}{independientes}, en el siguiente sentido combinatorio: para cada $n$,
el orden relativo de las $n$ primeras extracciones es uniforme entre
las $n!$ posibilidades, y $R_n = {}$``la $n$-ésima extracción es un
récord (mayor que todas las anteriores)''. Admitiendo que los
\omterm{def:b2:proba:space}{sucesos} $R_n$ son \omterm{def:b2:proba:independence}{independientes} con $\P(R_n) = 1/n$ (demuéstrese
al menos esta última igualdad por simetría), prueba usando
Borel--Cantelli que ocurren infinitos récords casi seguramente, pero
que los récords en instantes \emph{consecutivos} $n, n+1$ ocurren
infinitas veces con probabilidad: calcula $\sum_n \P(R_n \cap
R_{n+1})$ y concluye qué da Borel--Cantelli~1.
\end{exercise}

\begin{exercise}[$\star\star\star$]\label{exo:b2:proba:8}
(Al estilo de Kochen--Stone, versión más fácil) Sean $(A_n)$
\omterm{def:b2:proba:independence}{sucesos independientes} con $\P(A_n) = \frac{1}{n+1}$. Prueba que
$\P(\limsup A_n) = 1$, aunque $\P(A_n) \to 0$: ``individualmente
raros, colectivamente ciertos''. Recíprocamente, exhibe una sucesión
de \omterm{def:b2:proba:space}{sucesos} (dependientes) con $\sum\P(A_n) = \infty$ y $\P(\limsup
A_n) = 0$, lo que muestra que no puede prescindirse de la
\omterm{def:b2:proba:independence}{independencia} en Borel--Cantelli~2.
\end{exercise}

\begin{exercise}[$\star$]\label{exo:b2:proba:9}
Se lanza una moneda con probabilidad de cara $p \in \intoo01$ hasta
la primera cara. Calcula la probabilidad de que esto ocurra en un
rango impar, y evalúala para una moneda equilibrada.
\end{exercise}

\begin{exercise}[$\star\star$]\label{exo:b2:proba:10}
Sean $(A_n)_{n\geq1}$ \omterm{def:b2:proba:independence}{sucesos independientes} con $\P(A_n) = p_n < 1$.
Prueba que
\[
\P\Bigl(\bigcap_{n\geq1}A_n^c\Bigr)
= \prod_{n\geq1}(1 - p_n)
:= \lim_{N\to\infty}\prod_{n=1}^N(1 - p_n),
\]
y que ese límite es $> 0$ si y solo si $\sum p_n < \infty$.
Reconcílialo con Borel--Cantelli: cuando $\sum p_n = \infty$, no solo
ocurre algún $A_n$ casi seguramente, sino que ocurren infinitos.
\end{exercise}

\begin{exercise}[$\star\star$]\label{exo:b2:proba:11}
(Las cajas de cerillas de Banach) Un fumador lleva una caja de $n$
cerillas en cada bolsillo y mete la mano cada vez en un bolsillo
elegido uniformemente al azar. Cuando encuentra por primera vez una
caja vacía, ¿cuál es la probabilidad de que la otra caja contenga
exactamente $k$ cerillas? Prueba que la respuesta es
$\binom{2n-k}{n}2^{-(2n-k)}$ y comprueba que esas probabilidades
suman $1$ para $n = 1$.
\end{exercise}

\begin{exercise}[$\star\star\star$]\label{exo:b2:proba:12}
(La $\sigma$-aditividad es un axioma real) (a) Prueba que no hay
ninguna \omterm{def:b2:proba:space}{medida de probabilidad} sobre $(\N, \mathcal P(\N))$ que dé el
mismo peso a todos los conjuntos unitarios. (b) Para $A \subseteq
\N^*$, sea $d(A) = \lim_n\frac{\abs{A\cap\intint1n}}{n}$ cuando el
límite existe (la \emph{densidad natural}). Prueba que $d$ es
finitamente aditiva sobre las parejas en las que existen las tres
densidades, que da densidad $0$ a todo conjunto unitario y densidad
$1$ a $\N^*$, y concluye que $d$ no es $\sigma$-aditiva. (c) Exhibe
un conjunto sin densidad. \emph{(Altérnense bloques
$\intint{2^{2k}}{2^{2k+1}-1}$ dentro y fuera.)}
\end{exercise}

\section{Problema: el paseo aleatorio simple sobre $\Z$ es recurrente}

\begin{omfigure}
\begin{tikzpicture}[scale=0.42, y=0.9cm]
  \draw[->] (0,0) -- (25,0) node[below] {$n$};
  \draw[->] (0,-2.6) -- (0,2.9) node[left] {$S_n$};
  \draw[omDef, thick] plot coordinates {(0,0) (1,1) (2,0)
    (3,-1) (4,0) (5,1) (6,2) (7,1) (8,2) (9,1) (10,0) (11,-1)
    (12,-2) (13,-1) (14,0) (15,-1) (16,0) (17,1) (18,2) (19,1)
    (20,0) (21,1) (22,0) (23,1) (24,2)};
  \foreach \t in {2,4,10,14,16,20,22}
    \fill[omThm] (\t,0) circle (5pt);
  \fill (0,0) circle (5pt);
\end{tikzpicture}

{\small Veinticuatro pasos de un \omterm{pb:b2:proba:1}{paseo aleatorio simple}; los puntos
rojos marcan los retornos al origen. El problema demuestra que, con
probabilidad $1$, esos puntos no dejan nunca de aparecer; y sin
embargo, el tiempo de espera entre ellos tiene media divergente.}
\end{omfigure}

\begin{problem}[{Problema de fin de semana --- el teorema de
recurrencia de Pólya sobre $\Z$, con el problema de la papeleta y el
sabor del arcoseno por el camino}]\label{pb:b2:proba:1}
Lánzese una moneda equilibrada para siempre; sea $X_i = \pm1$ el
$i$-ésimo paso y $S_n = X_1 + \dots + X_n$ el \emph{paseo aleatorio
simple}\index{paseo aleatorio} sobre $\Z$, con $S_0 = 0$. Como en el~\cref{ex:b2:proba:runs}, todos los \omterm{def:b2:proba:space}{sucesos} de más abajo están
determinados por un número finito de lanzamientos o son
combinaciones \omterm{def:b2:structures:countable}{numerables} de tales \omterm{def:b2:proba:space}{sucesos}, y la \omterm{def:b2:proba:independence}{independencia} de
los \omterm{def:b2:proba:space}{sucesos} que dependen de bloques disjuntos de lanzamientos forma
parte del modelo. Escribimos $u_n = \P(S_{2n} = 0)$ y $N_n(k)$ para
el número de caminos de $\pm1$ de \omterm{def:b2:curves:length}{longitud} $n$ que van de $0$ a $k$.

\medskip
\noindent\textbf{Parte I --- Contar caminos.}
\begin{enumerate}
  \item Prueba que $N_n(k) = \binom{n}{(n+k)/2}$ cuando $n + k$ es
        par y $\abs k \leq n$, y $0$ en caso contrario; deduce que
        $\P(S_n = k) = N_n(k)\,2^{-n}$. ¿Por qué es igualmente
        probable cada camino individual de \omterm{def:b2:curves:length}{longitud} $n$?
  \item Prueba que $S_{2n+1} \neq 0$, que $u_n =
        \binom{2n}{n}4^{-n}$, y calcula $u_1, u_2, u_3$.
  \item Demuestra que $u_n = \frac{2n-1}{2n}\,u_{n-1}$; deduce que
        $(u_n)$ decrece a $0$ y, a partir del~\cref{ex:b2:comparison:centralbinomial}, que
        \[
        u_n \sim \frac{1}{\sqrt{\pi n}},
        \qquad\text{de modo que}\qquad
        \sum_n u_n = \infty .
        \]
  \item (Principio de reflexión) Para $k \geq 1$, prueba que los
        caminos de longitud $n$ de $1$ a $k$ que tocan $0$ están en
        biyección con los caminos de $-1$ a $k$; deduce que el
        número de caminos de $0$ a $k$ que permanecen $> 0$ después
        del instante $0$ es $N_{n-1}(k-1) - N_{n-1}(k+1)$.
  \item (Teorema de la papeleta) Deduce que
        \[
        \P\bigl(S_1 > 0, \dots, S_{n-1} > 0 \bigm| S_n =
        k\bigr) = \frac kn \qquad (k \geq 1) :
        \]
        en un recuento en el que el ganador aventaja por $k$ de $n$
        papeletas, la probabilidad de que el ganador fuera por
        delante durante todo el escrutinio es $k/n$. Verifícalo a
        mano para $n = 3$, $k = 1$.
\end{enumerate}

\medskip
\noindent\textbf{Parte II --- El retorno al origen.}
\begin{enumerate}[resume]
  \item Demuestra la identidad clave
        \[
        \P(S_1 \neq 0,\ S_2 \neq 0,\ \dots,\ S_{2n} \neq 0) =
        u_n
        \]
        \emph{(condiciónese al primer paso, súmense los recuentos de
        la pregunta 4 sobre el punto final y telescópese; remátese
        con $2\binom{2n-1}{n} = \binom{2n}{n}$)}.
  \item Deduce de la \omterm{def:b2:metric:continuity}{continuidad} monótona
        (\cref{thm:b2:proba:continuity}) que el paseo vuelve a $0$
        al menos una vez con probabilidad $1$, y que $f_n :=
        \P(\text{primer retorno en el instante }2n)$ cumple
        \[
        f_n = u_{n-1} - u_n = \frac{u_n}{2n-1},
        \qquad \sum_{n\geq1}f_n = 1 .
        \]
  \item Prueba que $\sum_n 2n\,f_n = \infty$: el retorno es cierto,
        pero la serie que calcularía el tiempo medio de espera
        diverge (en el vocabulario del~\cref{ch:b2:randomvar}, el
        tiempo de retorno tiene esperanza infinita).
  \item Demuestra que, para todo $k \geq 1$, $\P(\text{al menos }
        k\text{ retornos a }0) = 1$ \emph{(descompóngase sobre los
        instantes de los $k$ primeros retornos: los bloques de
        lanzamientos correspondientes son disjuntos, de modo que las
        probabilidades se multiplican y suman $(\sum_nf_n)^k$)}; y
        concluye con la \omterm{def:b2:metric:continuity}{continuidad} monótona:
        \[
        \P(S_n = 0 \text{ para infinitos } n) = 1 :
        \]
        el \omterm{pb:b2:proba:1}{paseo aleatorio simple} sobre $\Z$ es \emph{recurrente}.
  \item Prueba que el paseo visita todo sitio $k \in \Z$ casi
        seguramente y, por tanto (por recurrencia, reiniciando en la
        primera visita), infinitas veces. \emph{(Los signos de las
        excursiones sucesivas desde $0$ son monedas equilibradas
        \omterm{def:b2:proba:independence}{independientes}; una excursión positiva visita $1$.)}
\end{enumerate}

\medskip
\noindent\textbf{Parte III --- Borel--Cantelli y el paseo sesgado.}
\begin{enumerate}[resume]
  \item Los \omterm{def:b2:proba:space}{sucesos} $A_n = \{S_{2n} = 0\}$ cumplen $\sum\P(A_n) =
        \infty$; explica por qué Borel--Cantelli~2 \emph{no} se les
        aplica, y qué daría Borel--Cantelli~1 si la serie
        convergiera. (Esta es la estrategia de toda la parte.)
  \item Sea ahora la moneda de sesgo $p \neq \frac12$, con $q = 1 -
        p$. Prueba que $\P(S_{2n} = 0) = \binom{2n}n(pq)^n =
        u_n\,(4pq)^n$ con $4pq < 1$, deduce que $\sum_n\P(S_{2n} =
        0) < \infty$ y concluye por Borel--Cantelli~1 que el paseo
        sesgado vuelve a $0$ solo un número finito de veces, casi
        seguramente.
  \item Todavía para $p \neq \frac12$: prueba que $\P(S_n = k) \leq
        \binom{n}{\floor{n/2}}\,(pq)^{n/2}\,(p/q)^{k/2}$ para cada
        $k$ fijo, deduce que todo sitio se visita un número finito
        de veces casi seguramente y concluye que $\abs{S_n} \to
        \infty$ casi seguramente: el paseo sesgado es
        \emph{transitorio}.
  \item De vuelta a la moneda equilibrada: usando la pregunta 6,
        calcula la probabilidad de que $200$ lanzamientos no
        produzcan \emph{ningún} empate ($S_n \neq 0$ para $1 \leq n
        \leq 200$), numéricamente $u_{100} \approx 0.056$. Comenta
        el lento decaimiento $1/\sqrt{\pi n}$: los empates son
        ciertos a la larga, pero más raros de lo que sugiere la
        intuición.
  \item (Primer paso por un punto) Sea $T_1$ el primer instante en
        que el paseo alcanza $1$. Usando el principio de reflexión
        para el máximo $M_n = \max_{i\leq n}S_i$ (demostrado en la
        pregunta 16, que no depende de esta), o directamente a
        partir de la pregunta 7 condicionando al primer paso, prueba
        que $\P(T_1 = 2n - 1) = f_n$; deduce que $\P(T_1 < \infty) =
        1$, mientras que la serie del tiempo medio $\sum(2n-1)f_n$
        diverge.
\end{enumerate}

\medskip
\noindent\textbf{Parte IV --- Máximos, último cero, ventajas
prolongadas.}
\begin{enumerate}[resume]
  \item (Reflexión para el máximo) Para $k \geq 1$, demuestra que
        \[
        \P(M_n \geq k) = 2\,\P(S_n > k) + \P(S_n = k)
        \]
        reflejando el camino tras su primera visita al nivel $k$.
  \item Deduce que $\P(M_{2n} \geq 1) = 1 - u_n$, es decir, $\P(S_i
        \leq 0 \text{ para todo } i \leq 2n) = u_n$: la probabilidad
        de no ir nunca por delante coincide con la probabilidad de
        no estar nunca en cero (pregunta 6); dos \omterm{def:b2:proba:space}{sucesos} distintos,
        una misma probabilidad.
  \item (Último cero) Sea $L_{2n} = \max\{k \leq 2n : S_k = 0\}$
        (par). Combinando la pregunta 6 con la \omterm{def:b2:proba:independence}{independencia} de los
        bloques disjuntos de lanzamientos, prueba que
        \[
        \P(L_{2n} = 2k) = u_k\,u_{n-k}
        \qquad (0 \leq k \leq n),
        \]
        y deduce, sin más cálculo, la identidad binomial
        $\sum_{k=0}^n u_ku_{n-k} = 1$.
  \item Prueba que la ley de $L_{2n}$ es \omterm{def:b2:quadratic:adjoint}{simétrica} ($\P(L = 2k) =
        \P(L = 2n - 2k)$) y, usando $u_j \sim 1/\sqrt{\pi j}$, que
        sus extremos son sus valores más probables. Tabula para $n =
        5$: $\P(L_{10} = 0) = u_5 \approx 0.246$ frente a
        $\P(L_{10} = 4) = u_2u_3 \approx 0.117$. Interprétalo: en
        una partida equilibrada larga, el último empate tiende a ser
        muy temprano o muy tardío; las ventajas prolongadas son la
        regla, no la excepción.
  \item Ensambla las preguntas 16--19 en un párrafo sobre la imagen
        de las fluctuaciones del paseo equilibrado: la escala
        difusiva que sugiere la pregunta 3, la certeza del retorno
        frente al tiempo medio de espera divergente, y la
        persistencia de las ventajas con sabor a arcoseno.
\end{enumerate}

\medskip
\noindent\textbf{Parte V --- La identidad de renovación y el teorema
de Pólya.}
\begin{enumerate}[resume]
  \item Demuestra, partiendo $\{S_{2n} = 0\}$ según el instante del
        primer retorno, la \emph{identidad de renovación}
        \[
        u_n = \sum_{k=1}^n f_k\,u_{n-k} \quad (n \geq 1),
        \qquad\text{y por tanto}\qquad
        U(x)\bigl(1 - F(x)\bigr) = 1 \quad (0 \leq x < 1),
        \]
        donde $U(x) = \sum_{n\geq0}u_nx^n$ y $F(x) =
        \sum_{n\geq1}f_nx^n$ (justifica los radios y el producto de
        series con el~\cref{ch:b2:powerseries}).
  \item Deduce la \emph{dicotomía de recurrencia}: haciendo $x \to
        1^-$ (límites monótonos de series de coeficientes no
        negativos),
        \[
        \sum_n u_n = \infty \iff \sum_n f_n = 1 ,
        \]
        y compruébala contra las preguntas 3 y 7 (paseo equilibrado)
        y 12 (paseo sesgado).
  \item (Dimensión $2$) El paseo simple sobre $\Z^2$ da pasos
        $(\pm1, 0)$, $(0, \pm1)$ uniformemente. Prueba que las
        coordenadas giradas $U_n = X_n + Y_n$ y $V_n = X_n - Y_n$
        realizan paseos equilibrados \emph{\omterm{def:b2:proba:independence}{independientes}} sobre $\Z$,
        deduce que
        \[
        \P\bigl(S^{(2)}_{2n} = (0,0)\bigr) = u_n^2 \sim
        \frac1{\pi n},
        \qquad \sum_n u_n^2 = \infty ,
        \]
        y concluye con las preguntas 21--22 (cuyas demostraciones se
        transfieren palabra por palabra) que el paseo sobre $\Z^2$
        es recurrente.
  \item (Dimensión $3$) Para el paseo simple sobre $\Z^3$, admítase
        la estimación local $\P(S^{(3)}_{2n} = 0) \leq C\,n^{-3/2}$
        (demostrada con el teorema local del límite en el volumen
        del tercer año). Deduce de Borel--Cantelli~1 que el paseo
        sobre $\Z^3$ es transitorio, y enuncia el resultado
        \omterm{def:b2:metric:complete}{completo}: el \emph{teorema de Pólya}: el
        \omterm{pb:b2:proba:1}{paseo aleatorio simple} es recurrente en las dimensiones $1$ y
        $2$, y transitorio en dimensión $3$ y superiores.
  \item Síntesis. Enumera el papel exacto que desempeñan: el recuento
        de caminos y la reflexión; la \omterm{def:b2:metric:continuity}{continuidad} monótona; la
        \omterm{def:b2:proba:independence}{independencia} de los bloques disjuntos de lanzamientos;
        Borel--Cantelli~1; y la identidad de renovación. ¿Qué único
        hecho \omterm{def:b2:powerseries:analytic}{analítico} ($u_n \sim 1/\sqrt{\pi n}$, de donde $\sum
        u_n = \infty$, pero también $\sum u_n^2 = \infty$ y $\sum
        n^{-3/2} < \infty$) decide entre recurrencia y transitoriedad
        en cada dimensión?
\end{enumerate}
\end{problem}
